\documentclass[aspectratio=169,11pt]{beamer}

% Compile with XeLaTeX or LuaLaTeX on Overleaf
\usepackage{kotex}
\usepackage{amsmath, amssymb, bm, mathtools}
\usepackage{booktabs}
\usepackage{hyperref}
\usepackage{xcolor}

\usetheme{Madrid}
\usecolortheme{default}
\setbeamertemplate{navigation symbols}{}
\setbeamertemplate{footline}[frame number]

\title{Flow Matching}
\subtitle{정의, 초기조건, 시간에 따른 변화, 그리고 Conditional Flow Matching Loss}
\author{}
\date{}

\newcommand{\R}{\mathbb{R}}
\newcommand{\E}{\mathbb{E}}
\newcommand{\pt}{p_t}
\newcommand{\vx}{\mathbf{x}}
\newcommand{\vz}{\mathbf{z}}
\newcommand{\vu}{\mathbf{u}}
\newcommand{\vv}{\mathbf{v}}

\begin{document}

% Slide 1
\begin{frame}
  \titlepage
\end{frame}

% Slide 2
\begin{frame}{1. Flow Matching이란?}
\small
\textbf{핵심 아이디어:} 쉬운 분포 $p_0(\vx)$에서 시작해서, 목표 데이터 분포 $p_1(\vx)$로 가는
\alert{연속적인 흐름(velocity field)}을 학습한다.

\vspace{0.5em}
\begin{block}{상태의 진화}
\[
\frac{d\vx_t}{dt} = \vv_\theta(\vx_t, t), \qquad t \in [0,1].
\]
\end{block}

\begin{itemize}
  \item $\vx_0 \sim p_0$: 보통 가우시안 같은 쉬운 초기 분포
  \item $\vx_1 \sim p_1$: 우리가 생성하고 싶은 데이터 분포
  \item 목표는 $\vv_\theta$를 학습해서 ODE를 적분하면 $p_0 \to p_1$로 이동하게 만드는 것
\end{itemize}

\vspace{0.4em}
\textbf{직관적으로는:} noise를 아주 작은 시간 변화들을 통해 점점 data 쪽으로 밀어주는 벡터장을 배우는 것.
\end{frame}

% Slide 3
\begin{frame}{2. Initial Condition(초기조건)과 경로 설정}
\small
Flow Matching에서는 보통 두 endpoint를 둔다:
\[
\vx_0 \sim p_0, \qquad \vx_1 \sim p_1.
\]

그 다음 $(\vx_0, \vx_1)$를 잇는 \textbf{conditional probability path}를 정의한다. 가장 흔한 예시는 선형 경로:
\[
\vx_t = (1-t)\vx_0 + t\vx_1.
\]

보다 일반적으로는
\[
\vx_t = \alpha_t \vx_0 + \beta_t \vx_1,
\]
여기서
\[
\alpha_0 = 1,\; \beta_0 = 0, \qquad \alpha_1 = 0,\; \beta_1 = 1.
\]

\begin{block}{왜 Initial Condition이 중요한가?}
\begin{itemize}
  \item $t=0$에서 샘플은 반드시 $p_0$를 따라야 한다.
  \item $t=1$에서 샘플은 반드시 $p_1$를 따라야 한다.
  \item 즉, 경로의 시작과 끝을 고정해 두고 그 사이의 움직임을 학습하는 문제다.
\end{itemize}
\end{block}
\end{frame}

% Slide 4
\begin{frame}{3. 시간이 지나면서 어떻게 변하나?}
\small
조건부 경로가 주어지면, 각 시간 $t$에서의 움직임은 conditional velocity로 표현된다.

\[
\vu_t(\vx_t \mid \vx_0, \vx_1) = \frac{d\vx_t}{dt}.
\]

예를 들어 선형 경로 $\vx_t = (1-t)\vx_0 + t\vx_1$라면
\[
\vu_t(\vx_t \mid \vx_0, \vx_1) = \vx_1 - \vx_0.
\]

즉, 샘플은 시간에 따라 이 velocity field를 따라서 이동한다:
\[
\frac{d\vx_t}{dt} = \vu_t(\vx_t \mid \vx_0, \vx_1).
\]

이 흐름이 전체 분포 수준에서는 continuity equation을 만족한다:
\[
\partial_t p_t(\vx) + \nabla \cdot \bigl(p_t(\vx) \, \vv_t(\vx)\bigr) = 0.
\]

\begin{itemize}
  \item 개별 샘플은 ODE를 따라 움직이고
  \item 전체 분포는 시간이 지나며 $p_0 \to p_t \to p_1$로 변형된다.
\end{itemize}
\end{frame}

% Slide 5
\begin{frame}{4. Conditional Flow Matching Loss}
\small
학습에서는 실제 conditional velocity $\vu_t$를 직접 계산할 수 있는 경로를 설계하고,
모델 $\vv_\theta(\vx,t,\mathbf{c})$가 이를 회귀하도록 만든다.

\vspace{0.3em}
\textbf{Conditional FM loss:}
\[
\mathcal{L}_{\mathrm{CFM}}(\theta)
=
\E_{\mathbf{c},\,\vx_0,\,\vx_1,\,t}
\Bigl[
\| \vv_\theta(\vx_t,t,\mathbf{c}) - \vu_t(\vx_t\mid \vx_0,\vx_1,\mathbf{c}) \|^2
\Bigr].
\]

여기서
\begin{itemize}
  \item $\mathbf{c}$: 클래스, 텍스트, 조건 정보 등
  \item $\vx_1 \sim p_1(\vx \mid \mathbf{c})$: 조건부 데이터 분포
  \item $\vx_t$: $(\vx_0, \vx_1, \mathbf{c})$로부터 만든 중간 시점 샘플
\end{itemize}

\begin{block}{의미}
조건 $\mathbf{c}$가 주어졌을 때, 모델은 ``현재 위치 $\vx_t$에서 다음 순간 어디로 가야 하는지''를 예측한다.
즉, unconditional flow가 아니라 \alert{조건부 목표 분포로 향하는 flow}를 배운다.
\end{block}
\end{frame}

% Slide 6
\begin{frame}{5. Conditional Flow로 생성하기}
\small
학습이 끝나면, 원하는 조건 $\mathbf{c}$를 고정하고 초기 샘플에서 ODE를 적분한다:
\[
\frac{d\vx_t}{dt} = \vv_\theta(\vx_t,t,\mathbf{c}),
\qquad
\vx_0 \sim p_0.
\]

그러면 $t=1$에서 얻는 $\vx_1$은 목표 조건부 분포를 따르게 된다:
\[
\vx_1 \sim p_1(\vx \mid \mathbf{c}).
\]

\begin{block}{정리}
\begin{enumerate}
  \item 쉬운 초기 분포 $p_0$에서 시작한다.
  \item endpoint를 잇는 conditional path $\vx_t$를 만든다.
  \item 그 경로의 true velocity $\vu_t$를 계산한다.
  \item 모델 $\vv_\theta$가 이를 따라가도록 CFM loss로 학습한다.
  \item 생성 시에는 learned ODE를 따라 적분해서 조건부 샘플을 얻는다.
\end{enumerate}
\end{block}

\vspace{0.2em}
\textbf{한 줄 요약:} Flow Matching은 ``분포를 직접 맞추는 대신, 목표 분포로 이동하는 속도장''을 학습하는 방법이다.
\end{frame}

\end{document}
